\scanpage{24}\mspage{23}
% @ opening and heading
\deleted{\unsure{§|S}} En utilisant \uline{Conf} 3, \uline{Cont} 4, on trouve le

\uline{Scholie} Soit $F$ \inserted{$\subset M$} une figure. \deleted{L'application}
\[
\deleted{X\longmapsto\unsure{\mathcal{L}|\mathcal{L}_X}}
\]

% @ a), parts and partition
a) Les \deleted{ens \unsure{des}} \inserted{parties de $\mathcal{L}$} (``strates ouvertes de $F$'') $\mathcal{L}_X^{\circ}$ ($X\in F$) sont non vides et mutuellement disjointes, elles forment \deleted{de} une partition de $\mathcal{L}_{\deleted{X}F}=\bigcup_{X\in F}\mathcal{L}_X=\bigcup_{X\in F}\mathcal{L}_X^{\circ}$.

\deleted{b) L'applic} A fortiori, l'application
\[
 X\longmapsto \mathcal{L}_X^{\circ}\qquad F\longrightarrow\mathcal P(\mathcal{L})
\]
(injection)
% @ oblique caveat
% slanted
\margin{Mais bien sûr, si on ne se restreint pas aux multistrates qui font partie d'une figure donnée, la connaissance des $\mathcal{L}_X^{\circ}$ ou $\mathcal{L}_X$, ni même celle des $(\mathcal{L}_X,\mathcal{L}_X^{\circ})$, ne détermine $X$ (avec sa \uline{multi}-structure, \unsure{i.e.|ou} \unsure{ses} \unsure{parties} \unsure{fermées|ouvertes})}

% @ b), two maps
b) L'application
\[
 X\longmapsto \mathcal{L}_X\qquad F\longrightarrow\mathcal P(\mathcal{L})
\]
est également injective. \inserted{[donc si $X,Y$ sont comp. et $\mathcal{L}_X=\mathcal{L}_Y$ on a $X=Y$)} On a de plus, pour $X,Y\in F$, équivalence entre : les conditions suivantes
% @ four conditions, right margin
\[
\begin{aligned}
 \text{(i)}\quad &X\leq Y\\
 \text{(ii)}\quad &\mathcal{L}_X\subset \mathcal{L}_Y\\
 \text{(iii)}\quad &\mathcal{L}_X^{\circ}\subset \mathcal{L}_Y\\
 \text{(iv)}\quad &\mathcal{L}_X^{\circ}\cap \mathcal{L}_Y\ne\varnothing
\end{aligned}
\]
\margin{donc ces équivalences valent dès que $X,Y$ sont des multistrates \uline{compatibles}.}

% @ c) and right margin
c) Si $X,Y\in\deleted{\unsure{X}}F$, on a
\[
 \mathcal{L}_X\cap \mathcal{L}_Y=\bigcup_{\substack{Z\in F\\ Z\leq X,\,Z\leq Y}}\mathcal{L}_Z
\]
\margin{Itou pour c)}

% @ proof at foot
Dém de b). Les implications (i)$\Rightarrow$(ii)$\Rightarrow$(iii)$\Rightarrow$(iv) triviales, prouvons (iv)$\Rightarrow$(i) Comme $\mathcal{L}_Y=\bigcup_{\unsure{Z_i|Y_i}\leq Y} \mathcal{L}_{Y_i}^{\circ}$, \deleted{\unsure{on a}} $\exists Y_i$ avec $\mathcal{L}_X^{\circ}\cap \mathcal{L}_{Y_i}^{\circ}\ne\varnothing$ donc par a) on a $X=Y_i$ donc $X\leq Y$, qcd
